% slide-ready summary for the lecture box (max 6 lines); generated by the notebook
\begin{itemize}\setlength{\itemsep}{2pt}
  \item \textbf{Prețul ridicat al energiei în România.} Ianuarie--octombrie 2026: PZU 137,4 EUR/MWh, cu 23,4 peste Germania, dar $-$3,5 față de Ungaria --- o primă regională, nu națională.
  \item Prima față de Germania apare mai ales seara: 42,9 EUR/MWh între orele 18 și 22, față de 10,5 între 10 și 16; România exportă la prînz și importă seara.
  \item Descompunere contabilă în eșantion (fără interpretare cauzală), 2024$\to$2026: +39,9 EUR/MWh, din care costul marginal pe termen scurt al centralelor pe gaze, inclusiv carbonul (SRMC), +14,2, nuclear și hidro +37,2.
  \item Prețuri cu amănuntul (Eurostat, 2026-S1): gospodăriile plătesc în euro aproximativ media UE, dar cu 65\,\% peste ea în standarde ale puterii de cumpărare (PPS); pentru firme, componentele din 2025 plasează diferența în energie și furnizare.
  \item Pentru companii: acoperiți profilul de seară, evaluați stocarea și flexibilitatea, includeți în contractele de achiziție pe termen lung (PPA) clauze privind profilul orar.
\end{itemize}
